\documentclass[11pt,a4paper]{article}
% =====================================================================
%  Preámbulo común para los apuntes de Física Electrónica
%  Uso (desde unidades/XX-nombre/archivo.tex):
%     \documentclass[11pt,a4paper]{article}
%     % =====================================================================
%  Preámbulo común para los apuntes de Física Electrónica
%  Uso (desde unidades/XX-nombre/archivo.tex):
%     \documentclass[11pt,a4paper]{article}
%     % =====================================================================
%  Preámbulo común para los apuntes de Física Electrónica
%  Uso (desde unidades/XX-nombre/archivo.tex):
%     \documentclass[11pt,a4paper]{article}
%     \input{../../comun/preambulo.tex}
% =====================================================================

% ---------- Idioma y tipografía ----------
\usepackage[T1]{fontenc}
\usepackage[utf8]{inputenc}
\usepackage[spanish,es-nodecimaldot,es-noshorthands,es-tabla]{babel}
\usepackage{lmodern}
\usepackage{microtype}
\usepackage[a4paper,margin=2.3cm,headheight=14pt]{geometry}
\usepackage{parskip}
\setlength{\parskip}{0.55em}

% ---------- Matemática y unidades ----------
\usepackage{amsmath,amssymb,mathtools}
\usepackage{bm}
\usepackage[per-mode=symbol,output-decimal-marker={.},separate-uncertainty]{siunitx}
\sisetup{range-phrase={\text{--}},range-units=single}
\DeclareSIUnit{\eV}{eV}

% ---------- Gráficos ----------
\usepackage{graphicx}
\usepackage{xcolor}
\usepackage{tikz}
\usetikzlibrary{arrows.meta,decorations.pathmorphing,decorations.markings,
  patterns,calc,positioning,shapes.geometric,fit,backgrounds,shadings}
\usepackage{pgfplots}
\pgfplotsset{compat=1.18}
\usepgfplotslibrary{fillbetween}
\usepackage[siunitx,RPvoltages]{circuitikz}

% ---------- Tablas, listas, figuras ----------
\usepackage{booktabs,array,tabularx,multirow}
\usepackage{enumitem}
\setlist{itemsep=0.2em,topsep=0.3em}
\usepackage{caption,subcaption}
\captionsetup{font=small,labelfont=bf,format=hang}
\usepackage{float}
\usepackage{wrapfig}
\usepackage{needspace}

% ---------- Colores ----------
\definecolor{azul}{HTML}{1F4E79}
\definecolor{azulclaro}{HTML}{DCE9F5}
\definecolor{verde}{HTML}{2E7D32}
\definecolor{verdeclaro}{HTML}{E3F2E1}
\definecolor{naranja}{HTML}{C55A11}
\definecolor{naranjaclaro}{HTML}{FBE9DB}
\definecolor{rojo}{HTML}{B71C1C}
\definecolor{rojoclaro}{HTML}{FDECEA}
\definecolor{violeta}{HTML}{5E35B1}
\definecolor{violetaclaro}{HTML}{EDE7F6}
\definecolor{gris}{HTML}{555555}
% colores para dibujos de bandas / portadores
\colorlet{colEc}{azul}
\colorlet{colEv}{rojo}
\colorlet{colElec}{azul}
\colorlet{colHueco}{rojo}
\colorlet{colFoton}{naranja}

% ---------- Cajas didácticas ----------
\usepackage[most]{tcolorbox}
\tcbset{
  base/.style={enhanced,breakable,arc=2pt,boxrule=0.6pt,left=6pt,right=6pt,
    top=4pt,bottom=4pt,fonttitle=\bfseries\small,before skip=8pt,after skip=8pt}
}
% Definición
\newtcolorbox{definicion}[1][]{base,colback=azulclaro,colframe=azul,
  title={Definición\ifx\relax#1\relax\else: #1\fi}}
% Resultado / fórmula clave
\newtcolorbox{clave}[1][]{base,colback=verdeclaro,colframe=verde,
  title={\ifx\relax#1\relax Resultado clave\else #1\fi}}
% Ejemplo resuelto
\newtcolorbox{ejemplo}[1][]{base,colback=white,colframe=violeta,
  coltitle=white,colbacktitle=violeta,
  title={Ejemplo\ifx\relax#1\relax\else: #1\fi}}
% Corrección / aclaración respecto de las diapositivas
\newtcolorbox{ojo}[1][]{base,colback=rojoclaro,colframe=rojo,
  title={\ensuremath{\triangle}\ Ojo\ifx\relax#1\relax\else: #1\fi}}
% Intuición física / comentario
\newtcolorbox{intuicion}[1][]{base,colback=naranjaclaro,colframe=naranja,
  title={Idea física\ifx\relax#1\relax\else: #1\fi}}

% ---------- Encabezados ----------
\usepackage{fancyhdr}
\pagestyle{fancy}
\fancyhf{}
\renewcommand{\headrulewidth}{0.4pt}
\fancyhead[L]{\small\color{gris}Física Electrónica}
\fancyhead[R]{\small\color{gris}\unidadtitulo}
\fancyfoot[C]{\small\thepage}
\newcommand{\unidadtitulo}{}
\newcommand{\unidad}[1]{\renewcommand{\unidadtitulo}{#1}}

\usepackage{titlesec}
\titleformat{\section}{\Large\bfseries\color{azul}}{\thesection}{0.6em}{}
\titleformat{\subsection}{\large\bfseries\color{azul}}{\thesubsection}{0.6em}{}
\titleformat{\subsubsection}{\normalsize\bfseries\color{azul!80!black}}{\thesubsubsection}{0.6em}{}

\usepackage[hidelinks,bookmarksnumbered]{hyperref}
\usepackage{bookmark}

% ---------- Macros físicas ----------
\newcommand{\EG}{E_G}
\newcommand{\Ec}{E_c}
\newcommand{\Ev}{E_v}
\newcommand{\EF}{E_F}
\newcommand{\Eph}{E_{\mathrm{ph}}}
\newcommand{\lG}{\lambda_G}
\newcommand{\VT}{V_T}
\newcommand{\Popt}{P_{\mathrm{opt}}}
\newcommand{\Resp}{\mathcal{R}}
\newcommand{\vsat}{v_{\mathrm{sat}}}
\newcommand{\Voc}{V_{oc}}
\newcommand{\Isc}{I_{sc}}
\newcommand{\diapo}[1]{{\small\textsf{[diapo~#1]}}}
\newcommand{\dd}{\,\mathrm{d}}

% Distintivo para resultados que ya están en la hoja de fórmulas oficial.
% Uso: en el título de una caja, \begin{clave}[Título \enhoja], o en texto con \enhojatexto.
\newcommand{\enhojatexto}{{\scriptsize\sffamily\colorbox{verdeclaro}{\textcolor{verde!70!black}{\ensuremath{\checkmark}\,en la hoja}}}}
\newcommand{\enhoja}{\hfill{\scriptsize\sffamily\colorbox{white}{\textcolor{verde!70!black}{\ensuremath{\checkmark}\,en la hoja}}}}

\input{../../comun/constantes.tex}

% ---------- Estilos TikZ comunes ----------
\tikzset{
  banda/.style={line width=1.1pt},
  bandac/.style={banda,colEc},
  bandav/.style={banda,colEv},
  fermi/.style={dashed,gris,line width=0.6pt},
  elec/.style={circle,fill=colElec,inner sep=0pt,minimum size=5pt},
  hueco/.style={circle,draw=colHueco,fill=white,line width=0.9pt,inner sep=0pt,minimum size=5pt},
  foton/.style={colFoton,line width=0.9pt,decorate,
    decoration={snake,amplitude=1.6pt,segment length=5pt,post length=3pt},-{Stealth[length=5pt]}},
  flecha/.style={-{Stealth[length=5pt]},line width=0.7pt},
}

% =====================================================================

% ---------- Idioma y tipografía ----------
\usepackage[T1]{fontenc}
\usepackage[utf8]{inputenc}
\usepackage[spanish,es-nodecimaldot,es-noshorthands,es-tabla]{babel}
\usepackage{lmodern}
\usepackage{microtype}
\usepackage[a4paper,margin=2.3cm,headheight=14pt]{geometry}
\usepackage{parskip}
\setlength{\parskip}{0.55em}

% ---------- Matemática y unidades ----------
\usepackage{amsmath,amssymb,mathtools}
\usepackage{bm}
\usepackage[per-mode=symbol,output-decimal-marker={.},separate-uncertainty]{siunitx}
\sisetup{range-phrase={\text{--}},range-units=single}
\DeclareSIUnit{\eV}{eV}

% ---------- Gráficos ----------
\usepackage{graphicx}
\usepackage{xcolor}
\usepackage{tikz}
\usetikzlibrary{arrows.meta,decorations.pathmorphing,decorations.markings,
  patterns,calc,positioning,shapes.geometric,fit,backgrounds,shadings}
\usepackage{pgfplots}
\pgfplotsset{compat=1.18}
\usepgfplotslibrary{fillbetween}
\usepackage[siunitx,RPvoltages]{circuitikz}

% ---------- Tablas, listas, figuras ----------
\usepackage{booktabs,array,tabularx,multirow}
\usepackage{enumitem}
\setlist{itemsep=0.2em,topsep=0.3em}
\usepackage{caption,subcaption}
\captionsetup{font=small,labelfont=bf,format=hang}
\usepackage{float}
\usepackage{wrapfig}
\usepackage{needspace}

% ---------- Colores ----------
\definecolor{azul}{HTML}{1F4E79}
\definecolor{azulclaro}{HTML}{DCE9F5}
\definecolor{verde}{HTML}{2E7D32}
\definecolor{verdeclaro}{HTML}{E3F2E1}
\definecolor{naranja}{HTML}{C55A11}
\definecolor{naranjaclaro}{HTML}{FBE9DB}
\definecolor{rojo}{HTML}{B71C1C}
\definecolor{rojoclaro}{HTML}{FDECEA}
\definecolor{violeta}{HTML}{5E35B1}
\definecolor{violetaclaro}{HTML}{EDE7F6}
\definecolor{gris}{HTML}{555555}
% colores para dibujos de bandas / portadores
\colorlet{colEc}{azul}
\colorlet{colEv}{rojo}
\colorlet{colElec}{azul}
\colorlet{colHueco}{rojo}
\colorlet{colFoton}{naranja}

% ---------- Cajas didácticas ----------
\usepackage[most]{tcolorbox}
\tcbset{
  base/.style={enhanced,breakable,arc=2pt,boxrule=0.6pt,left=6pt,right=6pt,
    top=4pt,bottom=4pt,fonttitle=\bfseries\small,before skip=8pt,after skip=8pt}
}
% Definición
\newtcolorbox{definicion}[1][]{base,colback=azulclaro,colframe=azul,
  title={Definición\ifx\relax#1\relax\else: #1\fi}}
% Resultado / fórmula clave
\newtcolorbox{clave}[1][]{base,colback=verdeclaro,colframe=verde,
  title={\ifx\relax#1\relax Resultado clave\else #1\fi}}
% Ejemplo resuelto
\newtcolorbox{ejemplo}[1][]{base,colback=white,colframe=violeta,
  coltitle=white,colbacktitle=violeta,
  title={Ejemplo\ifx\relax#1\relax\else: #1\fi}}
% Corrección / aclaración respecto de las diapositivas
\newtcolorbox{ojo}[1][]{base,colback=rojoclaro,colframe=rojo,
  title={\ensuremath{\triangle}\ Ojo\ifx\relax#1\relax\else: #1\fi}}
% Intuición física / comentario
\newtcolorbox{intuicion}[1][]{base,colback=naranjaclaro,colframe=naranja,
  title={Idea física\ifx\relax#1\relax\else: #1\fi}}

% ---------- Encabezados ----------
\usepackage{fancyhdr}
\pagestyle{fancy}
\fancyhf{}
\renewcommand{\headrulewidth}{0.4pt}
\fancyhead[L]{\small\color{gris}Física Electrónica}
\fancyhead[R]{\small\color{gris}\unidadtitulo}
\fancyfoot[C]{\small\thepage}
\newcommand{\unidadtitulo}{}
\newcommand{\unidad}[1]{\renewcommand{\unidadtitulo}{#1}}

\usepackage{titlesec}
\titleformat{\section}{\Large\bfseries\color{azul}}{\thesection}{0.6em}{}
\titleformat{\subsection}{\large\bfseries\color{azul}}{\thesubsection}{0.6em}{}
\titleformat{\subsubsection}{\normalsize\bfseries\color{azul!80!black}}{\thesubsubsection}{0.6em}{}

\usepackage[hidelinks,bookmarksnumbered]{hyperref}
\usepackage{bookmark}

% ---------- Macros físicas ----------
\newcommand{\EG}{E_G}
\newcommand{\Ec}{E_c}
\newcommand{\Ev}{E_v}
\newcommand{\EF}{E_F}
\newcommand{\Eph}{E_{\mathrm{ph}}}
\newcommand{\lG}{\lambda_G}
\newcommand{\VT}{V_T}
\newcommand{\Popt}{P_{\mathrm{opt}}}
\newcommand{\Resp}{\mathcal{R}}
\newcommand{\vsat}{v_{\mathrm{sat}}}
\newcommand{\Voc}{V_{oc}}
\newcommand{\Isc}{I_{sc}}
\newcommand{\diapo}[1]{{\small\textsf{[diapo~#1]}}}
\newcommand{\dd}{\,\mathrm{d}}

% Distintivo para resultados que ya están en la hoja de fórmulas oficial.
% Uso: en el título de una caja, \begin{clave}[Título \enhoja], o en texto con \enhojatexto.
\newcommand{\enhojatexto}{{\scriptsize\sffamily\colorbox{verdeclaro}{\textcolor{verde!70!black}{\ensuremath{\checkmark}\,en la hoja}}}}
\newcommand{\enhoja}{\hfill{\scriptsize\sffamily\colorbox{white}{\textcolor{verde!70!black}{\ensuremath{\checkmark}\,en la hoja}}}}

% =====================================================================
%  Constantes oficiales de la materia ("Tabla de constantes.pdf").
%  Usar SIEMPRE estos valores en los ejemplos numéricos, para que
%  coincidan con los que se usan en clase y en el examen.
% =====================================================================

% ---------- Constantes físicas ----------
\newcommand{\qval}{\num{1.60e-19}}          % C
\newcommand{\kval}{\num{8.61e-5}}           % eV/K
\newcommand{\kJval}{\num{1.38e-23}}         % J/K
\newcommand{\hval}{\num{6.62e-34}}          % J s
\newcommand{\heVval}{\num{4.13e-15}}        % eV s
\newcommand{\epsval}{\num{8.85e-14}}        % F/cm
\newcommand{\VTval}{\SI{26}{mV}}            % kT/q a 300 K (0.0259 V)
\newcommand{\epsrox}{3.9}                   % SiO2

% ---------- Silicio ----------
\newcommand{\epsrSi}{11.8}
\newcommand{\niSi}{\num{1.0e10}}            % cm^-3
\newcommand{\munSi}{1350}                   % cm^2/(V s)
\newcommand{\mupSi}{480}                    % cm^2/(V s)
\newcommand{\EGSi}{1.12}                    % eV
\newcommand{\NCSi}{\num{3.22e19}}           % cm^-3
\newcommand{\NVSi}{\num{1.83e19}}           % cm^-3

% ---------- Germanio ----------
\newcommand{\epsrGe}{15.8}
\newcommand{\niGe}{\num{2.4e13}}
\newcommand{\munGe}{3900}
\newcommand{\mupGe}{1900}
\newcommand{\EGGe}{0.66}

% ---------- GaAs ----------
\newcommand{\epsrGaAs}{13.1}
\newcommand{\niGaAs}{\num{2.1e6}}
\newcommand{\munGaAs}{8800}
\newcommand{\mupGaAs}{400}
\newcommand{\EGGaAs}{1.42}

% ---------- Tabla para incluir en los apuntes ----------
\newcommand{\tablaconstantes}{%
\begin{tabular}{lccc}
\toprule
 & Si & Ge & GaAs \\
\midrule
$\epsilon_r$                     & 11.8 & 15.8 & 13.1 \\
$\mu_n$ (\si{cm^2/V.s})          & 1350 & 3900 & 8800 \\
$\mu_p$ (\si{cm^2/V.s})          & 480  & 1900 & 400 \\
$n_i$ (\si{cm^{-3}}, 300\,K)     & \num{1.0e10} & \num{2.4e13} & \num{2.1e6} \\
$E_G$ (\si{\eV})                 & 1.12 & 0.66 & 1.42 \\
\bottomrule
\end{tabular}}


% ---------- Estilos TikZ comunes ----------
\tikzset{
  banda/.style={line width=1.1pt},
  bandac/.style={banda,colEc},
  bandav/.style={banda,colEv},
  fermi/.style={dashed,gris,line width=0.6pt},
  elec/.style={circle,fill=colElec,inner sep=0pt,minimum size=5pt},
  hueco/.style={circle,draw=colHueco,fill=white,line width=0.9pt,inner sep=0pt,minimum size=5pt},
  foton/.style={colFoton,line width=0.9pt,decorate,
    decoration={snake,amplitude=1.6pt,segment length=5pt,post length=3pt},-{Stealth[length=5pt]}},
  flecha/.style={-{Stealth[length=5pt]},line width=0.7pt},
}

% =====================================================================

% ---------- Idioma y tipografía ----------
\usepackage[T1]{fontenc}
\usepackage[utf8]{inputenc}
\usepackage[spanish,es-nodecimaldot,es-noshorthands,es-tabla]{babel}
\usepackage{lmodern}
\usepackage{microtype}
\usepackage[a4paper,margin=2.3cm,headheight=14pt]{geometry}
\usepackage{parskip}
\setlength{\parskip}{0.55em}

% ---------- Matemática y unidades ----------
\usepackage{amsmath,amssymb,mathtools}
\usepackage{bm}
\usepackage[per-mode=symbol,output-decimal-marker={.},separate-uncertainty]{siunitx}
\sisetup{range-phrase={\text{--}},range-units=single}
\DeclareSIUnit{\eV}{eV}

% ---------- Gráficos ----------
\usepackage{graphicx}
\usepackage{xcolor}
\usepackage{tikz}
\usetikzlibrary{arrows.meta,decorations.pathmorphing,decorations.markings,
  patterns,calc,positioning,shapes.geometric,fit,backgrounds,shadings}
\usepackage{pgfplots}
\pgfplotsset{compat=1.18}
\usepgfplotslibrary{fillbetween}
\usepackage[siunitx,RPvoltages]{circuitikz}

% ---------- Tablas, listas, figuras ----------
\usepackage{booktabs,array,tabularx,multirow}
\usepackage{enumitem}
\setlist{itemsep=0.2em,topsep=0.3em}
\usepackage{caption,subcaption}
\captionsetup{font=small,labelfont=bf,format=hang}
\usepackage{float}
\usepackage{wrapfig}
\usepackage{needspace}

% ---------- Colores ----------
\definecolor{azul}{HTML}{1F4E79}
\definecolor{azulclaro}{HTML}{DCE9F5}
\definecolor{verde}{HTML}{2E7D32}
\definecolor{verdeclaro}{HTML}{E3F2E1}
\definecolor{naranja}{HTML}{C55A11}
\definecolor{naranjaclaro}{HTML}{FBE9DB}
\definecolor{rojo}{HTML}{B71C1C}
\definecolor{rojoclaro}{HTML}{FDECEA}
\definecolor{violeta}{HTML}{5E35B1}
\definecolor{violetaclaro}{HTML}{EDE7F6}
\definecolor{gris}{HTML}{555555}
% colores para dibujos de bandas / portadores
\colorlet{colEc}{azul}
\colorlet{colEv}{rojo}
\colorlet{colElec}{azul}
\colorlet{colHueco}{rojo}
\colorlet{colFoton}{naranja}

% ---------- Cajas didácticas ----------
\usepackage[most]{tcolorbox}
\tcbset{
  base/.style={enhanced,breakable,arc=2pt,boxrule=0.6pt,left=6pt,right=6pt,
    top=4pt,bottom=4pt,fonttitle=\bfseries\small,before skip=8pt,after skip=8pt}
}
% Definición
\newtcolorbox{definicion}[1][]{base,colback=azulclaro,colframe=azul,
  title={Definición\ifx\relax#1\relax\else: #1\fi}}
% Resultado / fórmula clave
\newtcolorbox{clave}[1][]{base,colback=verdeclaro,colframe=verde,
  title={\ifx\relax#1\relax Resultado clave\else #1\fi}}
% Ejemplo resuelto
\newtcolorbox{ejemplo}[1][]{base,colback=white,colframe=violeta,
  coltitle=white,colbacktitle=violeta,
  title={Ejemplo\ifx\relax#1\relax\else: #1\fi}}
% Corrección / aclaración respecto de las diapositivas
\newtcolorbox{ojo}[1][]{base,colback=rojoclaro,colframe=rojo,
  title={\ensuremath{\triangle}\ Ojo\ifx\relax#1\relax\else: #1\fi}}
% Intuición física / comentario
\newtcolorbox{intuicion}[1][]{base,colback=naranjaclaro,colframe=naranja,
  title={Idea física\ifx\relax#1\relax\else: #1\fi}}

% ---------- Encabezados ----------
\usepackage{fancyhdr}
\pagestyle{fancy}
\fancyhf{}
\renewcommand{\headrulewidth}{0.4pt}
\fancyhead[L]{\small\color{gris}Física Electrónica}
\fancyhead[R]{\small\color{gris}\unidadtitulo}
\fancyfoot[C]{\small\thepage}
\newcommand{\unidadtitulo}{}
\newcommand{\unidad}[1]{\renewcommand{\unidadtitulo}{#1}}

\usepackage{titlesec}
\titleformat{\section}{\Large\bfseries\color{azul}}{\thesection}{0.6em}{}
\titleformat{\subsection}{\large\bfseries\color{azul}}{\thesubsection}{0.6em}{}
\titleformat{\subsubsection}{\normalsize\bfseries\color{azul!80!black}}{\thesubsubsection}{0.6em}{}

\usepackage[hidelinks,bookmarksnumbered]{hyperref}
\usepackage{bookmark}

% ---------- Macros físicas ----------
\newcommand{\EG}{E_G}
\newcommand{\Ec}{E_c}
\newcommand{\Ev}{E_v}
\newcommand{\EF}{E_F}
\newcommand{\Eph}{E_{\mathrm{ph}}}
\newcommand{\lG}{\lambda_G}
\newcommand{\VT}{V_T}
\newcommand{\Popt}{P_{\mathrm{opt}}}
\newcommand{\Resp}{\mathcal{R}}
\newcommand{\vsat}{v_{\mathrm{sat}}}
\newcommand{\Voc}{V_{oc}}
\newcommand{\Isc}{I_{sc}}
\newcommand{\diapo}[1]{{\small\textsf{[diapo~#1]}}}
\newcommand{\dd}{\,\mathrm{d}}

% Distintivo para resultados que ya están en la hoja de fórmulas oficial.
% Uso: en el título de una caja, \begin{clave}[Título \enhoja], o en texto con \enhojatexto.
\newcommand{\enhojatexto}{{\scriptsize\sffamily\colorbox{verdeclaro}{\textcolor{verde!70!black}{\ensuremath{\checkmark}\,en la hoja}}}}
\newcommand{\enhoja}{\hfill{\scriptsize\sffamily\colorbox{white}{\textcolor{verde!70!black}{\ensuremath{\checkmark}\,en la hoja}}}}

% =====================================================================
%  Constantes oficiales de la materia ("Tabla de constantes.pdf").
%  Usar SIEMPRE estos valores en los ejemplos numéricos, para que
%  coincidan con los que se usan en clase y en el examen.
% =====================================================================

% ---------- Constantes físicas ----------
\newcommand{\qval}{\num{1.60e-19}}          % C
\newcommand{\kval}{\num{8.61e-5}}           % eV/K
\newcommand{\kJval}{\num{1.38e-23}}         % J/K
\newcommand{\hval}{\num{6.62e-34}}          % J s
\newcommand{\heVval}{\num{4.13e-15}}        % eV s
\newcommand{\epsval}{\num{8.85e-14}}        % F/cm
\newcommand{\VTval}{\SI{26}{mV}}            % kT/q a 300 K (0.0259 V)
\newcommand{\epsrox}{3.9}                   % SiO2

% ---------- Silicio ----------
\newcommand{\epsrSi}{11.8}
\newcommand{\niSi}{\num{1.0e10}}            % cm^-3
\newcommand{\munSi}{1350}                   % cm^2/(V s)
\newcommand{\mupSi}{480}                    % cm^2/(V s)
\newcommand{\EGSi}{1.12}                    % eV
\newcommand{\NCSi}{\num{3.22e19}}           % cm^-3
\newcommand{\NVSi}{\num{1.83e19}}           % cm^-3

% ---------- Germanio ----------
\newcommand{\epsrGe}{15.8}
\newcommand{\niGe}{\num{2.4e13}}
\newcommand{\munGe}{3900}
\newcommand{\mupGe}{1900}
\newcommand{\EGGe}{0.66}

% ---------- GaAs ----------
\newcommand{\epsrGaAs}{13.1}
\newcommand{\niGaAs}{\num{2.1e6}}
\newcommand{\munGaAs}{8800}
\newcommand{\mupGaAs}{400}
\newcommand{\EGGaAs}{1.42}

% ---------- Tabla para incluir en los apuntes ----------
\newcommand{\tablaconstantes}{%
\begin{tabular}{lccc}
\toprule
 & Si & Ge & GaAs \\
\midrule
$\epsilon_r$                     & 11.8 & 15.8 & 13.1 \\
$\mu_n$ (\si{cm^2/V.s})          & 1350 & 3900 & 8800 \\
$\mu_p$ (\si{cm^2/V.s})          & 480  & 1900 & 400 \\
$n_i$ (\si{cm^{-3}}, 300\,K)     & \num{1.0e10} & \num{2.4e13} & \num{2.1e6} \\
$E_G$ (\si{\eV})                 & 1.12 & 0.66 & 1.42 \\
\bottomrule
\end{tabular}}


% ---------- Estilos TikZ comunes ----------
\tikzset{
  banda/.style={line width=1.1pt},
  bandac/.style={banda,colEc},
  bandav/.style={banda,colEv},
  fermi/.style={dashed,gris,line width=0.6pt},
  elec/.style={circle,fill=colElec,inner sep=0pt,minimum size=5pt},
  hueco/.style={circle,draw=colHueco,fill=white,line width=0.9pt,inner sep=0pt,minimum size=5pt},
  foton/.style={colFoton,line width=0.9pt,decorate,
    decoration={snake,amplitude=1.6pt,segment length=5pt,post length=3pt},-{Stealth[length=5pt]}},
  flecha/.style={-{Stealth[length=5pt]},line width=0.7pt},
}

\unidad{Guía 5 -- BJT: resoluciones}

\newcommand{\aF}{\alpha_F}
\newcommand{\aR}{\alpha_R}
\newcommand{\aT}{\alpha_T}
\newcommand{\eVEB}{\left(e^{V_{EB}/V_T}-1\right)}
\newcommand{\eVCB}{\left(e^{V_{CB}/V_T}-1\right)}
\newcommand{\eVBE}{\left(e^{V_{BE}/V_T}-1\right)}
\newcommand{\eVBC}{\left(e^{V_{BC}/V_T}-1\right)}
\newcommand{\guia}{\textsf{[resp.\ guía]}}
\newcommand{\prob}[1]{\needspace{6\baselineskip}\section*{Problema #1}\addcontentsline{toc}{section}{Problema #1}}

\title{\color{azul}\textbf{Guía 5 -- Transistor bipolar de juntura (BJT)}\\[0.3em]
  \large Resoluciones completas y resultados verificados}
\author{Física Electrónica (23.09, ITBA)}
\date{Complemento del resumen de las unidades 7 y 8}

\begin{document}
\maketitle
\thispagestyle{empty}

\begin{tcolorbox}[base,colback=white,colframe=gris,title={Cómo usar este solucionario}]
\small
\begin{itemize}
  \item Cada problema tiene: planteo, desarrollo y \textbf{resultado en caja verde}. Las cajas
  \textcolor{rojo}{\textbf{Ojo}} marcan diferencias con las respuestas al final de la guía
  (\guia) y confusiones típicas; están todas resumidas en la tabla de correcciones del final.
  \item Todos los números se recalcularon con un script (Python). Constantes: tabla de la
  materia ($q=\qval$\,C, $k=\kval$\,eV/K, $\epsilon_0=\epsval$\,F/cm, $\epsilon_r=\epsrSi$,
  $\mu_n=\munSi$, $\mu_p=\mupSi$\,\si{cm^2/V.s}); $n_i$ el que da el enunciado
  ($\num{1.5e10}\,\si{cm^{-3}}$).
  \item $V_T=kT/q=\SI{25.83}{mV}$ a \SI{300}{K}. Con \SI{26}{mV} los resultados cambian unos pocos
  por ciento (las exponenciales son sensibles); las conclusiones no cambian.
  \item Convención de signos y notación: la de la hoja de fórmulas (PNP, $I_{F0}$, $\aF$, $\gamma$,
  $\aT$). Para NPN se invierten tensiones y se intercambian electrones$\leftrightarrow$huecos.
\end{itemize}
\end{tcolorbox}

\subsection*{Resultados rápidos}
\begin{table}[H]
\centering\small
\begin{tabular}{@{}lp{11.5cm}@{}}
\toprule
Prob. & Resultado \\
\midrule
1 (PNP) & $V_{EB}=\SI{0.648}{V}$, $V_{CB}=\SI{-2}{V}$: $W=\SI{0.918}{\micro m}$, $I_{Ep}=\SI{0.488}{mA}$,
$I_{En}=\SI{167}{\micro A}$, $I_{RB}=\SI{16.1}{\micro A}$, $I_B=\SI{183}{\micro A}$,
$\alpha=0.721$, $\beta=2.58$, $\gamma=0.745$ (tabla completa en el problema). \\
2 (NPN) & mismo punto: $I_{En}=\SI{1.354}{mA}$, $I_B=\SI{116}{\micro A}$, $\alpha=0.920$, $\beta=11.6$, $\gamma=0.931$. \\
3 y 4 & $x_{EB}=\SI{158}{nm}$, $x_{CB}=\SI{85}{nm}$, $W_{\text{ef}}=\SI{1.757}{\micro m}$ (igual para PNP y NPN). \\
5 & $I_{CB0}=(1-\aF\aR)I_{R0}$; además $I_{CE0}=I_{CB0}/(1-\aF)=(\beta_F+1)I_{CB0}$. \\
6 & $g_m=\SI{387}{mA/V}$ (\SI{385}{mA/V} con \SI{26}{mV}), $r_\pi\approx\SI{1.29}{k\ohm}$ (\SI{1.30}{k\ohm}), $r_o=\SI{80}{k\ohm}$. \\
7 & $h_{ie}=r_x+r_\pi$, $h_{re}\approx r_\pi/r_\mu\approx0$, $h_{fe}=g_mr_\pi=\beta_0$, $h_{oe}=1/r_o\,(+\beta/r_\mu)$. \\
8 & Los tres cortocircuitos dan una ecuación de diodo $I=I_0\left(e^{V/V_T}-1\right)$. \\
9 & NPN; electrones. (b) y (c): $\alpha$ y $\beta$ disminuyen. \\
10 & a) F \quad b) F \quad c) F \quad d) F \quad e) F. \\
14 & $V_{BE}=V_T\ln\!\left(1+I_1/I_{F0}\right)$. \\
15 & $H=-g_m\,(r_o\parallel R_L)\,\dfrac{r_\pi}{r_\pi+R_{in}}$. \\
\bottomrule
\end{tabular}
\end{table}

\tableofcontents
\newpage

% =====================================================================
\prob{1 -- PNP: componentes de corriente, $\alpha$, $\beta$, $\gamma$}
% =====================================================================

\textbf{Datos.} $N_{AE}=\num{1e16}$, $N_{DB}=\num{8e15}$, $N_{AC}=\SI{1e13}{cm^{-3}}$,
$W_B=\SI{1}{\micro m}$, $A=\SI{1e-5}{cm^2}$, $\tau_p=\tau_n=\SI{10}{ns}$, $T=\SI{300}{K}$,
$n_i=\SI{1.5e10}{cm^{-3}}$.

\subsection*{Paso 1: parámetros de cada región}
En cada zona importan los \textbf{minoritarios}: electrones en el emisor y colector (p), huecos en
la base (n).
\[
  D=\mu V_T,\qquad L=\sqrt{D\tau}
\]
\begin{table}[H]
\centering\small
\begin{tabular}{lcccc}
\toprule
Región & minoritario & $D$ (\si{cm^2/s}) & $L$ (\si{\micro m}) & concentración de equilibrio (\si{cm^{-3}}) \\
\midrule
Emisor (p) & e$^-$ & $1350\cdot0.02583=34.87$ & 5.91 & $n_{E0}=n_i^2/N_{AE}=\num{22500}$ \\
Base (n) & h$^+$ & $480\cdot0.02583=12.40$ & 3.52 & $p_{B0}=n_i^2/N_{DB}=\num{28125}$ \\
Colector (p) & e$^-$ & 34.87 & 5.91 & $n_{C0}=n_i^2/N_{AC}=\num{2.25e7}$ \\
\bottomrule
\end{tabular}
\end{table}

\subsection*{Paso 2: ancho efectivo de la base}
Potenciales de contacto:
\[
  V_{bi,EB}=V_T\ln\frac{N_{AE}N_{DB}}{n_i^2}=\SI{0.687}{V},\qquad
  V_{bi,CB}=V_T\ln\frac{N_{AC}N_{DB}}{n_i^2}=\SI{0.509}{V}.
\]
Sólo interesa la parte de cada zona desierta que \emph{invade la base} (lado n):
\[
  x_{n}=\sqrt{\frac{2\epsilon_s}{q}\,\frac{N_A}{N_D\,(N_A+N_D)}\,(V_{bi}-V_A)}
\]
con $V_A=V_{EB}$ para la juntura E--B y $V_A=V_{CB}$ para la C--B. Entonces
$W=W_B-x_{EB}-x_{CB}$.

\begin{table}[H]
\centering\small
\begin{tabular}{cccccc}
\toprule
$V_{EB}$ (V) & $V_{CB}$ (V) & $V_{bi,EB}-V_{EB}$ (mV) & $x_{EB}$ (nm) & $x_{CB}$ (nm) & $W$ (\si{\micro m}) \\
\midrule
0.628 & $-2$ & 59 & 73.1 & 22.6 & 0.9043 \\
0.648 & $-2$ & 39 & 59.5 & 22.6 & 0.9179 \\
0.668 & $-2$ & 19 & 41.5 & 22.6 & 0.9359 \\
0.648 & $-10$ & 39 & 59.5 & 46.3 & 0.8943 \\
0.648 & $-20$ & 39 & 59.5 & 64.6 & 0.8759 \\
\bottomrule
\end{tabular}
\end{table}
Coinciden con la guía. Como $N_{AC}\ll N_{DB}$, la zona desierta C--B está casi toda en el
colector (\SI{18}{\micro m} del lado p para \SI{-2}{V}) y en la base entra muy poco: por eso el
efecto Early de este transistor es moderado.

\subsection*{Paso 3: componentes de corriente}
\begin{clave}[Corrientes del BJT PNP (base de ancho arbitrario) \enhoja]
\[
  I_{Ep}=qA\frac{D_B}{L_B}p_{B0}\left[\frac{e^{V_{EB}/V_T}-1}{\tanh(W/L_B)}-\frac{e^{V_{CB}/V_T}-1}{\sinh(W/L_B)}\right]
\]
\[
  I_{Cp}=qA\frac{D_B}{L_B}p_{B0}\left[\frac{e^{V_{EB}/V_T}-1}{\sinh(W/L_B)}-\frac{e^{V_{CB}/V_T}-1}{\tanh(W/L_B)}\right]
\]
\[
  I_{En}=qA\frac{D_E}{L_E}n_{E0}\eVEB,\qquad
  I_{Cn}=-qA\frac{D_C}{L_C}n_{C0}\eVCB
\]
\[
  I_{B-R}=I_{Ep}-I_{Cp},\qquad I_B=I_E-I_C=I_{En}+I_{B-R}-I_{Cn}
\]
\[
  \alpha=\frac{I_C}{I_E},\quad \beta=\frac{I_C}{I_B},\quad \gamma=\frac{I_{Ep}}{I_E},\quad \aT=\frac{I_{Cp}}{I_{Ep}}
\]
\end{clave}
(La hoja escribe $I_B=I_{En}+I_{Cn}+I_{RB}$; con el signo de $I_{Cn}$ definido arriba, $I_{Cn}>0$
en MAD y corresponde restarlo. Es del orden de pA, así que numéricamente no cambia nada.)

\begin{ejemplo}[cálculo detallado para $V_{EB}=\SI{0.648}{V}$, $V_{CB}=\SI{-2}{V}$]
$e^{V_{EB}/V_T}=e^{25.087}=\num{7.85e10}$, \quad $W/L_B=0.9179/3.521=0.2607$.
\[
  qA\frac{D_B}{L_B}p_{B0}=\num{1.6e-19}\cdot\num{1e-5}\cdot\frac{12.40}{\num{3.521e-4}}\cdot28125
  =\SI{1.585e-15}{A}
\]
\[
  I_{Ep}=\num{1.585e-15}\left[\frac{\num{7.85e10}}{\tanh 0.2607}+\frac{1}{\sinh 0.2607}\right]
  =\SI{0.488}{mA},\qquad I_{Cp}=\SI{0.472}{mA}
\]
\[
  I_{En}=\num{1.6e-24}\cdot\frac{34.87}{\num{5.905e-4}}\cdot22500\cdot\num{7.85e10}=\SI{167}{\micro A},
\]
\[
  I_{Cn}=\num{1.6e-24}\cdot\frac{34.87}{\num{5.905e-4}}\cdot\num{2.25e7}=\SI{2.13}{pA}
\]
\[
  I_{B-R}=\SI{16.1}{\micro A},\quad I_B=\SI{183.1}{\micro A},\quad
  \alpha=\frac{0.4721}{0.6553}=0.7205,\quad \beta=\frac{0.4721}{0.1831}=2.58,\quad \gamma=0.745
\]
\textbf{Control rápido} con las aproximaciones de base corta de la hoja:
\[
  \gamma\approx\frac{1}{1+\dfrac{D_E}{D_B}\dfrac{N_B}{N_E}\dfrac{W}{L_E}}
  =\frac{1}{1+2.8125\cdot0.8\cdot0.1554}=0.741,\qquad
  \aT\approx1-\frac{W^2}{2L_B^2}=0.966 .
\]
\end{ejemplo}

\begin{clave}[Resultados Problema 1 (PNP)]
\centering\small
\setlength{\tabcolsep}{3.5pt}
\begin{tabular}{cccccccccccc}
\toprule
$V_{EB}$ & $V_{CB}$ & $I_{Ep}$ & $I_{Cp}$ & $I_{En}$ & $I_{Cn}$ & $I_{B-R}$ & $I_B$ & $\alpha$ & $\beta$ & $\gamma$ & $\aT$\\
(V) & (V) & (mA) & (mA) & (\si{\micro A}) & (pA) & (\si{\micro A}) & (\si{\micro A}) & & & & \\
\midrule
0.628 & $-2$  & 0.2284 & 0.2210 & 77.0  & 2.13 & 7.33  & 84.3  & 0.7239 & 2.62 & 0.7479 & 0.9679\\
0.648 & $-2$  & 0.4883 & 0.4721 & 167.0 & 2.13 & 16.13 & 183.1 & 0.7205 & 2.58 & 0.7451 & 0.9670\\
0.668 & $-2$  & 1.0396 & 1.0040 & 362.2 & 2.13 & 35.67 & 397.9 & 0.7162 & 2.52 & 0.7416 & 0.9657\\
\midrule
0.648 & $-10$ & 0.5006 & 0.4849 & 167.0 & 2.13 & 15.72 & 182.7 & 0.7263 & 2.65 & 0.7499 & 0.9686\\
0.648 & $-20$ & 0.5107 & 0.4953 & 167.0 & 2.13 & 15.40 & 182.4 & 0.7308 & 2.72 & 0.7536 & 0.9698\\
\bottomrule
\end{tabular}
\end{clave}

\begin{ojo}[la tabla de respuestas de la guía tiene un error de unidades]
La guía da $I_{Ep}\approx\SI{46}{mA}$, $I_{B-R}\approx\SI{157}{nA}$, $\beta\approx286$ y
$\gamma\approx0.9965$. Esos valores se reproducen \emph{exactamente} (dentro de las diferencias
de $V_T$) si se usa $W$ en \si{\micro m}$\times10^{-6}$ (metros) mezclado con el resto en
\si{cm}: $I_{Ep}\propto1/W$ queda $\times100$ e $I_{B-R}\propto W$ queda $/100$.
Las columnas que no dependen de $W$ ($I_{En}$, $I_{Cn}$) sí coinciden con la guía.

Prueba física de que $\beta\approx286$ es imposible: con $N_E/N_B=1.25$ el emisor casi no está más
dopado que la base, y la fórmula de la hoja da $\gamma\approx0.74$, o sea
$\beta\le\gamma/(1-\gamma)\approx2.9$. Un $\beta$ de cientos necesita $N_E\gg N_B$.
\end{ojo}

\begin{ojo}[validez del modelo en estos puntos]
Los números son los del modelo ideal que pide el problema, pero el punto elegido está al borde
de su validez:
$V_{bi}-V_{EB}$ es apenas \SIrange{19}{59}{mV} (la aproximación de vaciamiento es dudosa) y
$p_B(0)=p_{B0}e^{V_{EB}/V_T}=\num{1.0e15}$, \num{2.2e15}, \SI{4.8e15}{cm^{-3}}, que es
\SIrange{13}{60}{\percent} de $N_{DB}$: ya no hay baja inyección para $V_{EB}=\SI{0.668}{V}$.
\end{ojo}

\subsection*{Conclusiones}
\begin{itemize}
  \item \textbf{$V_{EB}$ controla exponencialmente.} $+\SI{20}{mV}$ multiplica las corrientes por
  $\approx e^{0.02/0.02583}=2.17$ (en la tabla: $\times2.14$, porque además $W$ crece un poco).
  \item \textbf{$V_{EB}\uparrow\Rightarrow\beta\downarrow$ levemente:} $x_{EB}$ se achica, $W$
  crece, $I_{Cp}\propto1/W$ baja y $I_{B-R}\propto W$ sube.
  \item \textbf{$|V_{CB}|\uparrow\Rightarrow$ efecto Early:} $x_{CB}$ crece, $W$ baja, a $V_{EB}$
  fijo aumentan $I_{Ep}$ e $I_C$ (de \SI{0.472}{mA} a \SI{0.495}{mA}) y $\beta$ ($2.58\to2.72$),
  mientras baja la recombinación en la base. $I_{En}$ no cambia (no depende de $W$).
  \item \textbf{$\beta$ muy bajo} porque está limitado por $\gamma\approx0.75$ (emisor mal dopado),
  no por el transporte ($\aT\approx0.97$). $I_{Cn}$ (pA) es despreciable y no depende de $V_{EB}$.
\end{itemize}

% =====================================================================
\prob{2 -- El mismo transistor en NPN}
% =====================================================================

Ahora $N_{DE}=\num{1e16}$, $N_{AB}=\num{8e15}$, $N_{DC}=\SI{1e13}{cm^{-3}}$. Los módulos de
dopaje son los mismos, así que \textbf{$V_{bi}$, $x_{EB}$, $x_{CB}$ y $W$ son idénticos} al
Problema~1. Lo que cambia es qué portador es minoritario en cada zona:
\begin{table}[H]
\centering\small
\begin{tabular}{lcccc}
\toprule
Región & minoritario & $D$ (\si{cm^2/s}) & $L$ (\si{\micro m}) & concentración \\
\midrule
Emisor (n) & h$^+$ & 12.40 & 3.52 & $p_{E0}=\num{22500}$ \\
Base (p) & e$^-$ & 34.87 & 5.91 & $n_{B0}=\num{28125}$ \\
Colector (n) & h$^+$ & 12.40 & 3.52 & $p_{C0}=\num{2.25e7}$ \\
\bottomrule
\end{tabular}
\end{table}
Las fórmulas son las del Problema~1 intercambiando $p\leftrightarrow n$ y usando $V_{BE}$, $V_{BC}$.
Aquí la corriente \emph{útil} es de electrones ($I_{En}$, $I_{Cn}$) y la parásita del emisor es de
huecos ($I_{Ep}$). (La guía mantiene los nombres de columna del PNP: su ``$I_{Ep}$'' es la
corriente útil.)

\begin{clave}[Resultados Problema 2 (NPN)]
\centering\small
\setlength{\tabcolsep}{3.5pt}
\begin{tabular}{cccccccccccc}
\toprule
$V_{BE}$ & $V_{BC}$ & $I_{En}$ & $I_{Cn}$ & $I_{Ep}$ & $I_{Cp}$ & $I_{B-R}$ & $I_B$ & $\alpha$ & $\beta$ & $\gamma$ & $\aT$\\
(V) & (V) & (mA) & (mA) & (\si{\micro A}) & (pA) & (\si{\micro A}) & (\si{\micro A}) & & & & \\
\midrule
0.628 & $-2$  & 0.6334 & 0.6260 & 45.9  & 1.27 & 7.35  & 53.3  & 0.9216 & 11.75 & 0.9324 & 0.9884\\
0.648 & $-2$  & 1.3537 & 1.3375 & 99.6  & 1.27 & 16.19 & 115.8 & 0.9203 & 11.55 & 0.9315 & 0.9880\\
0.668 & $-2$  & 2.8809 & 2.8451 & 216.0 & 1.27 & 35.81 & 251.8 & 0.9187 & 11.30 & 0.9303 & 0.9876\\
\midrule
0.648 & $-10$ & 1.3890 & 1.3732 & 99.6  & 1.27 & 15.78 & 115.4 & 0.9225 & 11.90 & 0.9331 & 0.9886\\
0.648 & $-20$ & 1.4177 & 1.4022 & 99.6  & 1.27 & 15.45 & 115.0 & 0.9242 & 12.19 & 0.9344 & 0.9891\\
\bottomrule
\end{tabular}
\end{clave}

\begin{ojo}[misma corrección que en el Problema 1]
La guía da corriente útil $\approx\SI{130}{mA}$, $I_{B-R}\approx\SI{157}{nA}$ y $\beta\approx1350$:
otra vez el factor 100 de $W$. Las columnas de $I_{Ep}$ ($\approx\SI{97}{\micro A}$) e $I_{Cp}$
(\SI{1.27}{pA}) sí coinciden.
\end{ojo}

\subsection*{Conclusiones y diferencias con el PNP}
\begin{itemize}
  \item Valen las mismas tendencias del Problema~1 (exponencial en $V_{BE}$, efecto Early).
  \item La corriente útil se multiplica por $\mu_n/\mu_p=2.81$: los electrones difunden mejor.
  \item La retroinyección al emisor (huecos) baja: $\propto D/L\propto\sqrt{D}$, o sea
  $\times\sqrt{480/1350}=0.60$. Entonces $\gamma$ pasa de 0.745 a 0.931.
  \item $L_B$ crece de 3.5 a \SI{5.9}{\micro m}, así que $\aT$ sube de 0.967 a 0.988.
  \item $I_{B-R}$ \textbf{casi no cambia} ($\approx qAn_{B0}W/2\tau\cdot e^{V/V_T}$ no depende de $D$).
  \item Resultado: $\beta$ sube de $\approx2.6$ a $\approx11.6$ ($\times4.5$). Por eso los BJT
  rápidos y de alta ganancia son NPN. Igualmente sigue siendo un mal transistor por el emisor
  poco dopado.
\end{itemize}

% =====================================================================
\prob{3 -- Ancho efectivo de la base (PNP)}
% =====================================================================

\textbf{Datos.} $N_{AE}=\num{6e17}$, $N_{DB}=\num{1e16}$, $N_{AC}=\SI{1e14}{cm^{-3}}$,
$W_B=\SI{2}{\micro m}$, $V_{EB}=\SI{0.605}{V}$, $V_{CB}=\SI{-5}{V}$, $T=\SI{300.15}{K}$
($V_T=\SI{25.84}{mV}$). Se usa $n_i=\SI{1.5e10}{cm^{-3}}$ como en el Problema~1 (con eso
coincide la guía). $\tau$ no interviene en $W$.

\textbf{Juntura E--B.}
\[
  V_{bi,EB}=0.02584\ln\frac{\num{6e17}\cdot\num{1e16}}{\num{2.25e20}}=\SI{0.7989}{V}
  \;\Rightarrow\; V_{bi}-V_{EB}=\SI{0.194}{V}
\]
\[
  x_{EB}=\sqrt{\frac{2\cdot\num{1.044e-12}\cdot0.194}{\num{1.6e-19}}\cdot
  \frac{\num{6e17}}{\num{1e16}\cdot\num{6.1e17}}}=\SI{1.578e-5}{cm}=\SI{158}{nm}
\]
\textbf{Juntura C--B.}
\[
  V_{bi,CB}=0.02584\ln\frac{\num{1e14}\cdot\num{1e16}}{\num{2.25e20}}=\SI{0.574}{V}
  \;\Rightarrow\; V_{bi}-V_{CB}=\SI{5.574}{V}
\]
\[
  x_{CB}=\sqrt{\frac{2\cdot\num{1.044e-12}\cdot5.574}{\num{1.6e-19}}\cdot
  \frac{\num{1e14}}{\num{1e16}\cdot\num{1.01e16}}}=\SI{8.49e-6}{cm}=\SI{85}{nm}
\]

\begin{clave}[Resultado Problema 3]
\[
  W_{\text{ef}}=W_B-x_{EB}-x_{CB}=2-0.158-0.085=\SI{1.757}{\micro m}
\]
\end{clave}
Con $n_i=\num{1e10}$ de la tabla: $x_{EB}=\SI{166}{nm}$, $x_{CB}=\SI{85}{nm}$,
$W_{\text{ef}}=\SI{1.749}{\micro m}$. La diferencia es mínima.

Aunque la juntura E--B está en directa, invade la base \emph{más} que la C--B en inversa,
porque del lado C el colector es 100 veces menos dopado que la base: la zona desierta se va casi
toda al colector.

% =====================================================================
\prob{4 -- Ídem para NPN}
% =====================================================================
Los anchos de vaciamiento sólo dependen de los \emph{módulos} de los dopajes y de las tensiones
($V_{bi}$ y $x_n$ no distinguen si el lado más dopado es p o n). Por eso:
\begin{clave}[Resultado Problema 4]
\[
  x_{EB}=\SI{158}{nm},\qquad x_{CB}=\SI{85}{nm},\qquad W_{\text{ef}}=\SI{1.757}{\micro m}
\]
\end{clave}
Lo que \emph{sí} cambia entre PNP y NPN son las corrientes ($D$, $L$ de los minoritarios), no la
geometría.

% =====================================================================
\prob{5 -- $I_{CB0}$ a partir de Ebers--Moll (NPN)}
% =====================================================================

Ebers--Moll para NPN ($I_E$ saliente, $I_C$ e $I_B$ entrantes):
\[
  I_E=I_{F0}\eVBE-\aR I_{R0}\eVBC,
\]
\[
  I_C=\aF I_{F0}\eVBE-I_{R0}\eVBC,
\]
con reciprocidad $\aF I_{F0}=\aR I_{R0}$.

\textbf{Emisor abierto} ($I_E=0$):
\[
  I_{F0}\eVBE=\aR I_{R0}\eVBC .
\]
Reemplazando en $I_C$:
\[
  I_C=\aF\aR I_{R0}\eVBC-I_{R0}\eVBC=-(1-\aF\aR)\,I_{R0}\eVBC .
\]
Con la juntura B--C fuertemente en inversa ($V_{BC}<-5V_T$, $e^{V_{BC}/V_T}\to0$):
\begin{clave}[Resultado Problema 5]
\[
  I_{CB0}=(1-\aF\aR)\,I_{R0}
\]
\end{clave}
\textbf{Detalle extra:} la tensión que queda en el emisor abierto es
$e^{V_{BE}/V_T}-1=-\aR I_{R0}/I_{F0}=-\aF$, o sea $V_{BE}=V_T\ln(1-\aF)<0$ (unos
\SI{-0.1}{V}): la juntura E--B queda levemente en inversa.

\textbf{Qué es físicamente.} Escribiendo los coeficientes con
$a=qA\frac{D_B}{L_B}n_{B0}$, $e=qA\frac{D_E}{L_E}p_{E0}$, $c=qA\frac{D_C}{L_C}p_{C0}$ y $u=W/L_B$:
\[
  I_{CB0}=c+a\,\frac{e\coth u+a}{e+a\coth u}
  \;\xrightarrow[\;W\ll L_B\;]{\gamma\to1}\;
  \underbrace{qA\frac{D_C}{L_C}p_{C0}}_{I_{Cp}\ \text{(colector)}}+\underbrace{\frac{qAn_{B0}W}{\tau_B}}_{\text{generación en la base}} .
\]
Domina el primer término (los minoritarios del colector, que es poco dopado), que es el $I_{Cp}$
que menciona la guía. Por eso $I_{CB0}\approx I_{Cp}$ es una buena aproximación, no una igualdad.

\textbf{Complemento: $I_{CE0}$} (base abierta, $I_B=0\Rightarrow I_E=I_C$). Repitiendo el álgebra:
\[
  I_{CE0}=\frac{(1-\aF\aR)I_{R0}}{1-\aF}=\frac{I_{CB0}}{1-\aF}=(\beta_F+1)\,I_{CB0}.
\]
La fuga con base abierta es amplificada por el transistor.

% =====================================================================
\prob{6 -- Modelo de pequeña señal ($\pi$ híbrido)}
% =====================================================================

\begin{figure}[H]
\centering
\begin{circuitikz}[scale=0.95,transform shape]
  \draw (0,2) node[left]{B} to[short,o-*] (1.5,2) to[R,l_=$r_\pi$] (1.5,0) to[short,*-] (0,0) node[left]{E} ;
  \draw (0,0) to[short,o-] (1.5,0);
  \draw (1.5,0) -- (6,0) to[short,-o] (7,0) node[right]{E};
  \draw (4,2) to[cI,l=$g_mv_\pi$] (4,0);
  \draw (6,2) to[R,l=$r_o$] (6,0);
  \draw (4,2) -- (6,2) to[short,*-o] (7,2) node[right]{C};
  \node at (2.2,1.6) {$+$}; \node at (2.2,0.4) {$-$}; \node at (2.4,1) {$v_\pi$};
  \draw (4,0) node[circ]{}; \draw (6,0) node[circ]{};
\end{circuitikz}
\caption{Modelo $\pi$ híbrido simplificado de baja frecuencia (sin $r_x$ ni $r_\mu$).}
\end{figure}

\[
  g_m=\frac{I_C}{V_T}=\frac{\SI{10}{mA}}{\SI{25.83}{mV}}=\SI{387}{mA/V},\qquad
  r_\pi=\frac{\beta}{g_m}=\frac{500}{0.387}=\SI{1.29}{k\ohm},
\]
\[
  r_o=\frac{V_A}{I_C}=\frac{\SI{800}{V}}{\SI{10}{mA}}=\SI{80}{k\ohm}.
\]
\begin{clave}[Resultado Problema 6]
$g_m=\SI{387}{mA/V}$, $r_\pi=\SI{1.29}{k\ohm}$, $r_o=\SI{80}{k\ohm}$.
Con $V_T=\SI{26}{mV}$ (como la guía): $g_m=\SI{384.6}{mA/V}$, $r_\pi=\SI{1.30}{k\ohm}$.
\end{clave}
Si se conoce $V_{CE}$, la expresión más precisa es $r_o=(V_A+V_{CE})/I_C$; con $V_A=\SI{800}{V}$ la
diferencia es despreciable.

% =====================================================================
\prob{7 -- Parámetros $h$ en emisor común a partir de Giacoletto}
% =====================================================================

Modelo de Giacoletto completo: $r_x$ ($r_{bb'}$), $r_\pi$ ($r_{b'e}$), $r_\mu$ ($r_{b'c}$),
$g_mv_\pi$ y $r_o$ ($r_{ce}$). Definición de parámetros híbridos:
\[
  v_1=h_{ie}\,i_1+h_{re}\,v_2,\qquad i_2=h_{fe}\,i_1+h_{oe}\,v_2 .
\]
\begin{figure}[H]
\centering
\begin{circuitikz}[scale=0.9,transform shape]
  \draw (0,2) node[left]{B} to[short,o-] (0.3,2) to[R,l=$r_x$] (2.3,2) node[above]{B$'$};
  \draw (2.3,2) to[short,*-] (2.3,2) to[R,l_=$r_\pi$] (2.3,0);
  \draw (2.3,2) -- (2.3,3.2) to[R,l=$r_\mu$] (6.3,3.2) -- (6.3,2);
  \draw (4.5,2) to[cI,l=$g_mv_\pi$] (4.5,0);
  \draw (6.3,2) to[R,l=$r_o$] (6.3,0);
  \draw (4.5,2) -- (6.3,2) to[short,*-o] (7.3,2) node[right]{C};
  \draw (0,0) node[left]{E} to[short,o-] (6.3,0) to[short,-o] (7.3,0) node[right]{E};
  \draw (4.5,0) node[circ]{}; \draw (2.3,0) node[circ]{}; \draw (6.3,0) node[circ]{};
  \node at (3.0,1.6) {$+$}; \node at (3.0,0.4) {$-$}; \node at (3.2,1) {$v_\pi$};
\end{circuitikz}
\caption{Modelo de Giacoletto (baja frecuencia: sin $C_\pi$ ni $C_\mu$).}
\end{figure}

\begin{itemize}
  \item \textbf{$h_{ie}$} ($v_2=0$): $v_\pi=i_1\,(r_\pi\parallel r_\mu)$ $\Rightarrow$
  $h_{ie}=r_x+r_\pi\parallel r_\mu\approx r_x+r_\pi$.
  \item \textbf{$h_{re}$} ($i_1=0$): divisor $r_\mu$--$r_\pi$ $\Rightarrow$
  $h_{re}=\dfrac{r_\pi}{r_\pi+r_\mu}\approx\dfrac{r_\pi}{r_\mu}\approx0$.
  \item \textbf{$h_{fe}$} ($v_2=0$): $i_2=g_mv_\pi-v_\pi/r_\mu$ $\Rightarrow$
  $h_{fe}=\left(g_m-\tfrac{1}{r_\mu}\right)(r_\pi\parallel r_\mu)\approx g_mr_\pi=\beta_0$.
  \item \textbf{$h_{oe}$} ($i_1=0$): $i_2=\dfrac{v_2}{r_o}+g_mv_\pi+\dfrac{v_2-v_\pi}{r_\mu}$
  con $v_\pi=v_2\frac{r_\pi}{r_\pi+r_\mu}$ $\Rightarrow$
  $h_{oe}=\dfrac{1}{r_o}+\dfrac{1+g_mr_\pi}{r_\pi+r_\mu}\approx\dfrac{1}{r_o}+\dfrac{\beta_0}{r_\mu}$.
\end{itemize}
\begin{clave}[Resultado Problema 7]
Modelo completo: $h_{ie}=r_x+r_\pi\parallel r_\mu$, $h_{re}=\dfrac{r_\pi}{r_\pi+r_\mu}$,
$h_{fe}=\left(g_m-\tfrac{1}{r_\mu}\right)(r_\pi\parallel r_\mu)$, $h_{oe}=\dfrac{1}{r_o}+\dfrac{1+g_mr_\pi}{r_\pi+r_\mu}$.
Simplificado ($r_x\to0$, $r_\mu\to\infty$, como la guía):
\[
  h_{ie}=r_\pi,\qquad h_{re}=0,\qquad h_{fe}=g_mr_\pi=\beta,\qquad h_{oe}=\frac{1}{r_o}.
\]
\end{clave}
La guía omite $h_{re}$: en el modelo simplificado vale 0 y conviene decirlo explícitamente.

% =====================================================================
\prob{8 -- BJT con dos terminales en cortocircuito}
% =====================================================================

Ebers--Moll PNP: $I_E=I_{F0}\eVEB-\aR I_{R0}\eVCB$, $I_C=\aF I_{F0}\eVEB-I_{R0}\eVCB$, $I_B=I_E-I_C$.
Hay \textbf{tres} formas de cortocircuitar dos terminales:

\textbf{(i) C--B en corto ($V_{CB}=0$): diodo E--(B,C).}
\[
  I_E=I_{F0}\eVEB,\qquad I_C=\aF I_E,\qquad I_B=(1-\aF)I_E .
\]
Una fracción $\aF$ de los huecos inyectados cruza la base y vuelve por el colector, que está en
corto con la base. Es un diodo casi ideal (factor de idealidad $\approx1$ en un rango amplio): es
el ``transistor conectado como diodo'' de los espejos de corriente (Problema~14).
\textbf{Ruptura baja}: la del E--B (\SIrange{5}{7}{V} típicos), porque el emisor está muy dopado
y el campo máximo se alcanza con poca tensión.

\textbf{(ii) E--B en corto ($V_{EB}=0$): diodo C--(B,E).}
\[
  I_C=-I_{R0}\eVCB,\qquad I_E=-\aR I_{R0}\eVCB=\aR I_C .
\]
Ahora el colector es el que inyecta, y sólo una fracción $\aR$ (chica) llega al emisor.
\textbf{Ruptura alta}: la juntura C--B (colector poco dopado, zona desierta ancha).

\textbf{(iii) E--C en corto ($V_{EB}=V_{CB}=V$): diodo B--(E,C).} Las dos junturas en paralelo:
\[
  I_B=I_E-I_C=\big[(1-\aF)I_{F0}+(1-\aR)I_{R0}\big]\left(e^{V/V_T}-1\right).
\]
La ruptura en inversa la impone la juntura más débil, la E--B.

\begin{clave}[Resultado Problema 8]
En los tres casos la corriente es $I=I_0\left(e^{V/V_T}-1\right)$: un diodo, con
$I_0=I_{F0}$, $I_{R0}$ o $(1-\aF)I_{F0}+(1-\aR)I_{R0}$ respectivamente.
$BV$: alta sólo en el caso (ii); (i) y (iii) quedan limitados por la juntura E--B.
\end{clave}

\begin{ojo}[sobre la comparación de $I_{F0}$ e $I_{R0}$]
La guía dice ``si el transistor se construye de forma tal que $I_{F0}>I_{R0}$''. Por reciprocidad,
$\aF I_{F0}=\aR I_{R0}$, y como $\aR<\aF$ en cualquier transistor normal, siempre resulta
$I_{R0}>I_{F0}$ (el colector, poco dopado y de mayor área, tiene más corriente de saturación).
Lo que hace ``mejor'' al diodo E--B no es una corriente de saturación mayor, sino que su corriente
es casi toda de huecos que atraviesan la base (alta $\gamma$).
\end{ojo}

% =====================================================================
\prob{9 -- Identificar el transistor y analizar dopaje y ancho de base}
% =====================================================================

\textbf{a)} En el emisor $E_G=\SI{1.15}{eV}$ y $|E_C-E_F|=\SI{0.22}{eV}<E_G/2=\SI{0.575}{eV}$: el
nivel de Fermi está en la mitad superior de la banda prohibida, así que el emisor es \textbf{tipo n}
y el transistor es \textbf{NPN}. Con $V_{BE}=+\SI{0.647}{V}$ (directa) y
$V_{BC}=\SI{-3.59}{V}$ (inversa) está en \textbf{MAD}. Por la juntura C--B pasan
\textbf{electrones} inyectados por el emisor que atravesaron la base sin recombinarse
(empujados por el campo de la juntura en inversa).

Para (b) y (c) conviene la expresión de la hoja (base corta, NPN):
\[
  \frac{1}{\beta}\approx\underbrace{\frac{D_E}{D_B}\,\frac{N_B}{N_E}\,\frac{W}{L_E}}_{\text{retroinyección }(\gamma)}
  +\underbrace{\frac{W^2}{2L_B^2}}_{\text{recombinación }(\aT)},
  \qquad \alpha=\frac{\beta}{\beta+1}.
\]

\textbf{b) $N_B\times2$, $W$ constante.} El primer término se duplica: la base inyecta más huecos
hacia el emisor en relación a los electrones que recibe ($n_{B0}=n_i^2/N_B$ se reduce a la mitad).
$\gamma\downarrow$. En el modelo ideal $\aT$ no cambia; en la realidad $\mu_n$ y $\tau_n$ bajan con
el dopaje, así que $L_B$ baja y $\aT$ también cae algo. \textbf{$\alpha\downarrow$ y
$\beta\downarrow$} (si $\gamma$ es lo que limita, $\beta$ cae aproximadamente a la mitad).
Como efecto colateral bueno: menor efecto Early y menor resistencia de base.

\textbf{c) $W\uparrow$, $N_B$ constante.} Suben los dos términos: $\gamma\downarrow$ (la corriente
útil $\propto1/W$ mientras la retroinyección es fija) y $\aT\downarrow$ (más tiempo de tránsito,
más recombinación, $\propto W^2$). \textbf{$\alpha\downarrow$ y $\beta\downarrow$}.

\begin{ojo}[mecanismo en (b)]
La guía lo atribuye a que ``más portadores se recombinan en la base''. Al duplicar $N_B$, la
concentración inyectada en la base se \emph{reduce} a la mitad, y con ella la recombinación: el
cociente $I_{B-R}/I_C$ no cambia en el modelo ideal. Lo que empeora es la \textbf{eficiencia de
emisión} $\gamma$. La conclusión ($\alpha$ y $\beta$ bajan) es correcta.
\end{ojo}

% =====================================================================
\prob{10 -- Verdadero o falso}
% =====================================================================
\begin{enumerate}[label=\alph*)]
  \item \textbf{F.} $V_A\approx Q_B/C_{jC}\propto W$. Para una misma variación $\Delta x_{CB}$, la
  variación relativa $\Delta W/W$ es menor en la base de \SI{800}{nm}: \textbf{menos} efecto Early.
  (Que sea epitaxial no cambia el argumento.)
  \item \textbf{F.} La juntura C--B está en inversa con tensiones grandes y variables: su zona
  desierta cambia mucho. La E--B está en directa a $V_{BE}\approx$ cte, y su $x_{EB}$ casi no varía.
  \item \textbf{F.} La primera mitad es correcta (aumento de $I_C$ con $V_{CE}$), pero el mecanismo es
  la \textbf{modulación del ancho de la base}. ``Modulación del largo del canal'' es el
  efecto análogo en el MOSFET; el BJT no tiene canal.
  \item \textbf{F.} En saturación las dos junturas inyectan: el colector inyecta en sentido contrario
  ($I_C$ baja) y aumenta la recombinación en la base ($I_B$ sube). $h_{FE}=I_C/I_B=\beta_{\text{forzado}}<\beta_F$:
  el máximo $h_{FE}$ se obtiene en MAD.
  \item \textbf{F.} En las dos polarizaciones, en las zonas neutras la corriente es de
  \textbf{difusión de minoritarios}. En directa son portadores \emph{inyectados} (mayoritarios de un
  lado que cruzan la barrera y pasan a ser minoritarios en exceso del otro). En inversa son
  minoritarios generados térmicamente que difunden hasta la juntura y el campo los
  \emph{arrastra} a través de ella (extracción). La guía dice ``directa: difusión de mayoritarios;
  inversa: arrastre de minoritarios'', que es una simplificación. Lo correcto: lo que cruza la
  barrera en directa son mayoritarios, pero lo que determina la corriente es su difusión como
  minoritarios.
\end{enumerate}

% =====================================================================
\prob{11 -- Preguntas teóricas}
% =====================================================================
\textbf{a) ¿Es lo mismo que dos diodos enfrentados?} No. En el BJT la base es mucho más angosta
que la longitud de difusión ($W\ll L_B$): los portadores inyectados por el emisor llegan al
colector sin recombinarse (acoplamiento, las fuentes $\aF I_F$, $\aR I_R$ de Ebers--Moll). En dos
diodos discretos, el ``medio'' (zona neutra más contacto metálico) es muchísimo más largo que $L$:
todo lo inyectado se recombina y no hay efecto transistor. Por eso dos diodos enfrentados sólo
conducen la corriente de saturación inversa del que queda en inversa (o la de avalancha); no hay
amplificación.

\textbf{b) $V_{CE,\text{sat}}<\SI{1}{V}$ con las dos junturas en directa.} Las tensiones
de juntura se \emph{restan} en la malla:
\[
  V_{CE}=V_{BE}-V_{BC}\approx0.7-0.5\approx\SI{0.2}{V}\quad\text{(NPN)}.
\]
$V_{BC}<V_{BE}$ porque la juntura colectora tiene mayor $I_{R0}$ (más área, colector poco dopado)
y necesita menos tensión para la misma corriente. De Ebers--Moll:
\[
  V_{CE,\text{sat}}=V_T\ln\frac{1+(\beta_{\text{forz}}+1)/\beta_R}{1-\beta_{\text{forz}}/\beta_F},
  \qquad\beta_{\text{forz}}=\frac{I_C}{I_B}<\beta_F .
\]
Ejemplo: $\beta_F=100$, $\beta_R=1$, $\beta_{\text{forz}}=10$ da $V_{CE,\text{sat}}=\SI{67}{mV}$
ideal; con las resistencias serie del colector, \SIrange{0.1}{0.3}{V} en la práctica.

\textbf{c) ¿Por qué la estructura planar vertical?} El efecto transistor necesita
$W\ll L_B$ (\si{\micro m} o menos). No se puede hacer una capa tan delgada y además contactarla
lateralmente con metal. En el proceso planar se difunde la base dentro del colector y después el
emisor dentro de la base: el ancho de base queda definido por la \emph{diferencia de dos
profundidades de difusión} (controlable con mucha precisión), y el efecto transistor ocurre sólo
en el volumen vertical debajo del emisor. Ventajas: colector grande (disipación y alto $BV$),
emisor inyectando sólo hacia abajo. Costo: resistencia de base parásita $r_x$ (la corriente de base
recorre la base angosta en forma lateral), que no se puede reducir subiendo $N_B$ porque bajaría
$\gamma$ y $\beta$.

\textbf{d) Dopajes ideales.} $N_E\gg N_B\gg N_C$ (típico $10^{19}$--$10^{20}$, $10^{17}$ y $10^{15}$--\SI{e16}{cm^{-3}}):
\begin{itemize}
  \item $N_E\gg N_B$: $\gamma\to1$ (poca retroinyección).
  \item $N_B$ moderado: suficiente para evitar el \emph{punch-through} y reducir el efecto Early y
  $r_x$, pero sin arruinar $\gamma$.
  \item $N_C$ bajo: zona desierta C--B ancha hacia el colector (alto $BV_{CB}$, baja $C_\mu$, poco
  Early porque casi no invade la base).
\end{itemize}

% =====================================================================
\prob{12 -- Diagrama de bandas de un NPN en equilibrio}
% =====================================================================
Ejemplo con dopajes típicos: $N_{DE}=\num{1e18}$, $N_{AB}=\num{1e17}$, $N_{DC}=\SI{1e15}{cm^{-3}}$,
$n_i=\num{1e10}$:
\[
  E_F-E_i\big|_E=V_T\ln\tfrac{N_{DE}}{n_i}=\SI{0.48}{eV},\quad
  E_i-E_F\big|_B=\SI{0.42}{eV},\quad
  E_F-E_i\big|_C=\SI{0.30}{eV},
\]
\[
  qV_{bi,EB}=0.48+0.42=\SI{0.89}{eV},\qquad qV_{bi,BC}=0.42+0.30=\SI{0.71}{eV}.
\]
\begin{figure}[H]
\centering
\begin{tikzpicture}[yscale=2.6,xscale=1.15]
  % regiones
  \fill[azulclaro!50] (0,-1.25) rectangle (2.8,1.2);
  \fill[rojoclaro!60] (3.4,-1.25) rectangle (5.6,1.2);
  \fill[azulclaro!50] (7.0,-1.25) rectangle (10,1.2);
  \node[font=\small\bfseries] at (1.4,1.12) {Emisor n$^+$};
  \node[font=\small\bfseries] at (4.5,1.12) {Base p};
  \node[font=\small\bfseries] at (8.5,1.12) {Colector n};
  % Ec
  \draw[bandac] (0,0.084) -- (2.8,0.084) to[out=0,in=180] (3.4,0.976) -- (5.6,0.976)
       to[out=0,in=180] (7.0,0.263) -- (10,0.263) node[right]{$E_C$};
  % Ev
  \draw[bandav] (0,-1.036) -- (2.8,-1.036) to[out=0,in=180] (3.4,-0.144) -- (5.6,-0.144)
       to[out=0,in=180] (7.0,-0.857) -- (10,-0.857) node[right]{$E_V$};
  % Ei
  \draw[gris,dotted,thick] (0,-0.476) -- (2.8,-0.476) to[out=0,in=180] (3.4,0.416) -- (5.6,0.416)
       to[out=0,in=180] (7.0,-0.297) -- (10,-0.297) node[right]{$E_i$};
  % EF
  \draw[fermi] (0,0) -- (10,0) node[right]{$E_F$};
  % Vbi
  \draw[flecha,<->] (2.3,0.084) -- node[left,font=\small]{$qV_{bi,EB}$} (2.3,0.976);
  \draw[gris,thin] (2.2,0.976) -- (3.4,0.976);
  \draw[flecha,<->] (7.6,0.263) -- node[right,font=\small]{$qV_{bi,BC}$} (7.6,0.976);
  \draw[gris,thin] (5.6,0.976) -- (7.7,0.976);
  % EG
  \draw[flecha,<->] (4.5,-0.144) -- node[right,font=\small]{$E_G$} (4.5,0.976);
  % zonas desiertas
  \draw[gris,dashed] (2.8,-1.25) -- (2.8,1.05); \draw[gris,dashed] (3.4,-1.25) -- (3.4,1.05);
  \draw[gris,dashed] (5.6,-1.25) -- (5.6,1.05); \draw[gris,dashed] (7.0,-1.25) -- (7.0,1.05);
  \node[font=\scriptsize,gris] at (3.1,-1.17) {ZCE};
  \node[font=\scriptsize,gris] at (6.3,-1.17) {ZCE};
\end{tikzpicture}
\caption{NPN en equilibrio. $E_F$ es plano (no hay corriente). En el emisor $E_F$ está cerca de
$E_C$, en la base cerca de $E_V$ y en el colector también cerca de $E_C$ pero menos que en el
emisor. La barrera E--B es mayor que la B--C. La ZCE E--B está mayormente en la base (el emisor
está más dopado); la B--C, mayormente en el colector.}
\end{figure}
Elementos a indicar: $E_C$, $E_V$, $E_i$, $E_F$ constante, $E_G$ igual en las tres zonas,
$qV_{bi,EB}>qV_{bi,BC}$, y las zonas desiertas ubicadas del lado menos dopado.

% =====================================================================
\prob{13 -- El BJT como llave para un LED}
% =====================================================================
\textbf{S1 abierta:} no hay corriente de base, $I_B=0$, la juntura B--E queda sin polarizar:
\textbf{corte}. $I_C\approx I_{CE0}\approx0$, el LED está apagado y $V_C\approx V_{CC}-V_{LED}$
(el LED casi no conduce).

\textbf{S1 cerrada:} $R_B$ fija la corriente de base
\[
  I_B=\frac{V_{CC}-V_{BE}}{R_B}\approx\frac{V_{CC}-0.7}{R_B},
\]
y por efecto transistor circula $I_C$ por $R_C$ y el LED, que se enciende. Para un LED de
\SI{700}{nm}: $E_{ph}=hc/\lambda=\SI{1.77}{eV}$, así que $V_{LED}\gtrsim\SI{1.8}{V}$.
\[
  V_C=V_{CC}-V_{LED}-I_CR_C .
\]
\begin{itemize}
  \item Si $V_C>V_B\approx\SI{0.7}{V}$ (B--C en inversa): MAD, $I_C=\beta I_B$.
  \item Si $\beta I_B$ ``pide'' más corriente de la que la malla puede dar, $V_C$ cae por debajo
  de $V_B$, B--C entra en directa: \textbf{saturación}, $V_{CE}\approx\SIrange{0.1}{0.2}{V}$ e
  $I_C$ queda fijada por el circuito externo:
  \[
    I_{C,\text{sat}}=\frac{V_{CC}-V_{LED}-V_{CE,\text{sat}}}{R_C}.
  \]
\end{itemize}
\begin{clave}[Resultado Problema 13]
Debe operar como \textbf{llave: corte (abierta) / saturación (cerrada)}. En saturación
$P_{BJT}=V_{CE}I_C$ es mínima y la corriente del LED queda fijada por $R_C$ (no depende de
$\beta$, que varía mucho entre unidades y con la temperatura). Condición de diseño:
\[
  \beta_{\min}I_B>I_{C,\text{sat}}\;\Longrightarrow\;
  R_B<\beta_{\min}\,R_C\,\frac{V_{CC}-0.7}{V_{CC}-V_{LED}-V_{CE,\text{sat}}}
\]
(en la práctica, con un factor de sobreexcitación de 2 a 10).
\end{clave}

% =====================================================================
\prob{14 -- Tensión base-emisor del transistor conectado como diodo}
% =====================================================================
En el circuito el colector está unido a la base: $V_{BC}=0$. La fuente $I_1$ entra por el
nodo B--C y sale por el emisor, así que $I_E=I_1$. Con Ebers--Moll (NPN):
\[
  I_E=I_{F0}\eVBE-\aR I_{R0}\underbrace{\left(e^{0}-1\right)}_{0}=I_{F0}\eVBE=I_1 .
\]
\begin{clave}[Resultado Problema 14]
\[
  V_{BE}=V_T\ln\!\left(1+\frac{I_1}{I_{F0}}\right)\approx V_T\ln\frac{I_1}{I_{F0}}
\]
\end{clave}
La guía escribe $I_S$: aquí $I_S$ es $I_{F0}$, la corriente de saturación de la juntura E--B con
$V_{CB}=0$. Si se usa la $I_S$ de transporte ($I_C=I_Se^{V_{BE}/V_T}$, $I_S=\aF I_{F0}$), queda
$V_{BE}=V_T\ln(1+\aF I_1/I_S)$. La diferencia es $V_T\ln\aF$, menor a \SI{1}{mV}. Además,
$I_C=\aF I_1$ e $I_B=(1-\aF)I_1$. Es la referencia de un espejo de corriente.

% =====================================================================
\prob{15 -- Ganancia de tensión del emisor común}
% =====================================================================
Se pasivan $V_{BB}$ y $V_{CC}$ (cortocircuitos en AC) y se reemplaza el transistor por el modelo
$\pi$ híbrido:
\begin{figure}[H]
\centering
\begin{circuitikz}[scale=0.9,transform shape]
  \draw (0,0) to[sV,l=$v_{in}$] (0,2) to[R,l=$R_{in}$] (2.5,2) to[short,-*] (2.5,2);
  \draw (2.5,2) to[R,l_=$r_\pi$] (2.5,0);
  \draw (4.5,2) to[cI,l=$g_mv_\pi$] (4.5,0);
  \draw (6.3,2) to[R,l=$r_o$] (6.3,0);
  \draw (8.1,2) to[R,l=$R_L$] (8.1,0);
  \draw (4.5,2) -- (8.1,2) to[short,-o] (9,2);
  \draw (0,0) -- (8.1,0) to[short,-o] (9,0);
  \draw (9,2) to[open,v^=$v_{out}$] (9,0);
  \node at (3.1,1.6) {$+$}; \node at (3.1,0.4) {$-$}; \node at (3.3,1) {$v_\pi$};
  \foreach \x in {2.5,4.5,6.3} {\draw (\x,0) node[circ]{};}
  \foreach \x in {6.3} {\draw (\x,2) node[circ]{};}
\end{circuitikz}
\caption{Equivalente de pequeña señal del emisor común.}
\end{figure}
\[
  v_\pi=v_{in}\frac{r_\pi}{r_\pi+R_{in}}\quad\text{(divisor)},\qquad
  v_{out}=-g_mv_\pi\,(r_o\parallel R_L)\quad\text{(la fuente inyecta hacia E)}.
\]
\begin{clave}[Resultado Problema 15]
\[
  H=\frac{v_{out}}{v_{in}}=-g_m\,\frac{r_\pi}{r_\pi+R_{in}}\,\frac{r_oR_L}{r_o+R_L}
  =-\frac{\beta\,(r_o\parallel R_L)}{r_\pi+R_{in}}
\]
\end{clave}
El signo menos indica inversión de fase de \ang{180}. Si $R_{in}\ll r_\pi$ y $R_L\ll r_o$:
$H\approx-g_mR_L$.

% =====================================================================
\section*{Tabla de correcciones respecto de las respuestas de la guía}
\addcontentsline{toc}{section}{Tabla de correcciones}
% =====================================================================
\begin{table}[H]
\centering\small
\begin{tabularx}{\textwidth}{@{}lX X@{}}
\toprule
Prob. & Dice la guía & Correcto \\
\midrule
1 & $I_{Ep}\approx\SI{46}{mA}$, $I_{B-R}\approx\SI{157}{nA}$, $\beta\approx286$, $\gamma\approx0.997$ &
Error de unidades en $W$ (factor 100). $I_{Ep}=\SI{0.488}{mA}$, $I_{B-R}=\SI{16.1}{\micro A}$,
$\beta=2.58$, $\gamma=0.745$. $\beta$ de cientos es imposible con $N_E/N_B=1.25$. \\
2 & corriente útil $\approx\SI{131}{mA}$, $\beta\approx1350$ & Mismo error: $\SI{1.354}{mA}$,
$\beta=11.6$, $\gamma=0.931$. \\
1, 2 & --- & El punto de trabajo viola baja inyección ($p_B(0)$ hasta $0.6N_B$) y $V_{EB}\approx V_{bi}$. \\
5 & ``demostrar $I_{CB0}=I_{Cp}$'' & Es aproximado: $I_{CB0}=I_{Cp}+$ generación en la base $+\ldots$ \\
7 & falta $h_{re}$ & $h_{re}=r_\pi/(r_\pi+r_\mu)\approx0$; si hay $r_x$, $h_{ie}=r_x+r_\pi$. \\
8 & ``si $I_{F0}>I_{R0}$\ldots''; dos casos & Por reciprocidad, siempre $I_{R0}>I_{F0}$. Hay un
tercer caso (E--C en corto). \\
9b & ``más recombinación en la base'' & El mecanismo dominante es la caída de $\gamma$. \\
10e & ``directa: difusión de mayoritarios; inversa: arrastre de minoritarios'' & En ambos casos,
difusión de minoritarios en las zonas neutras (inyección o extracción). \\
11a & ``no pueden conducir salvo avalancha'' & Conducen la corriente de saturación inversa (muy chica). \\
14 & $I_S$ & Es $I_{F0}$ (con $I_S$ de transporte aparece $\aF$). \\
\bottomrule
\end{tabularx}
\end{table}

\end{document}
